\section{The supplied reference policies are not strategic baselines}
\label{app:pilot}

An earlier project phase, before the present no-new-evaluation scope, audited the six supplied reference policies and the 130 static scenarios against the unchanged engine and ran twenty baseline-only games with no model queries. We report its results because they bear on question~8 and question~C. This draft ran nothing new.

The twenty baseline games applied 19,494 decisions. Nineteen reached the 1,000-decision observation limit and were treated as censored, with no winner assigned; censoring at that horizon is not evidence that the games would never end. No policy initiated a trade, and the games contain only 44 construction actions. Two repeated configurations reproduced their event streams exactly. On the static scenarios, 75 of 780 policy--scenario choices and 11 of 130 supplied reference actions failed the engine's legality check. Fifty-five scenarios failed at least one static consistency screen, mostly through house-stock discrepancies, and 68 passed both the reference-legality and static-invariant screens. Scoring rubrics in 128 of 130 scenarios award points for keywords in the private note, and all six policies choose identical actions in 41 scenarios.

These policies and scenarios are useful infrastructure, but they are not competent negotiating or developing opponents, and passing a legality screen does not certify a scenario as a capability probe. Any comparison of the LLM cohort against reference behavior would first need adjudicated tasks and opponents that trade and build. The pilot games are a separate population and are never pooled with the twenty LLM games.
